\documentclass[10pt]{beamer}
% \documentclass[handout,10pt]{beamer}
\usepackage{my_style}

\title[Math/CS 375]{Lecture 09 - B}
\subtitle{\S 4.1 Newton Interpolation }

\begin{document}
% -------------------------------------------------
\begin{frame}
  \titlepage
\end{frame}
% -------------------------------------------------

\begin{frame}{Goals}
We will answer these questions:\\[3mm]
\begin{itemize}
    \item What is the benefit of nesting polynomials? 
    \item What are the Newton polynomial basis functions?
    \item What type of linear system is involved in Newton interpolation?
    \item What are the advantages of Newton interpolation? 
    \item How do you use the Tabular method for calculating interpolation by hand?
\end{itemize}
\end{frame}



%%%%%%%%%%%%%%%%%%%%%%%%%%%%%%%%%%%%%%%%%%%
\sectionframe{Nesting Polynomials}
%%%%%%%%%%%%%%%%%%%%%%%%%%%%%%%%%%%%%%%%%%%
\begin{frame}{Polynomial Operation Count}
\begin{itemize}
    \item Consider the polynomial   
    $$p_3(x) = 5 + 4x +7x^2 + 2x^3$$

    \item When implemented in Matlab, it is 
        $$p_3(x) = \texttt{5 + 4*x + 7*x*x + 2*x*x*x}$$

    \item 3 additions AND (1+2+3) multiplications = 9 FLOPs\\[3mm]

   
\end{itemize}
\end{frame}

%%%%%%%%%%%%%%%%%%%%%%%%%%%%%%%%%%%%%%%%%%%
\begin{frame}{In General}

\begin{itemize}
    \item A general polynomial has the form 
    $$
    p_n(x) =  a_0 + a_1x+ a_2x^2+\cdots +a_nx^n
    $$
    \item There are exactly $n$ additions 
    \vfill 
    \item Zeroth term has zero multiplications
    \vfill 
    \item The first term has one multiplication... 
    \vfill 
    \item The second term has two multiplication... 
    \vfill 
\end{itemize}
\begin{block}{Number of FLOPs}
    
    In general, evaluating $p_n(x)$ costs $n$ additions and $\sum_{i=1}^n i$ multiplications or 
    $$\text{Total} = n+\left(\frac{n(n+1)}{2}\right) \text{ FLOPs}$$
\end{block}
\end{frame}

%%%%%%%%%%%%%%%%%%%%%%%%%%%%%%%%%%%%%%%%%%%%%%%
\begin{frame}{Factoring}
\begin{itemize}
    \item We can rewrite this polynomial 
    \begin{align*}
    p_3(x) &= 5 + 4x +7x^2 + 2x^3\\
           &= 5 + x(4 +7x + 2x^2)\\
           &= 5 + x(4 +x(7 + 2x))
    \end{align*}

    \item When implemented in Matlab, it is 
        $$p_3(x) = \texttt{5 + x*(4 + x*(7 + 2*x))}$$

    \item 3 additions AND 3 multiplications = 6 FLOPs\\[3mm]

\end{itemize}
\end{frame}

%%%%%%%%%%%%%%%%%%%%%%%%%%%%%%%%%%%%%%%%%%%
\begin{frame}{Optimal polynomial representation}
\begin{itemize}
    \item Given a polynomial
\begin{equation*}
  p(x) = -5 + 4x -7x^2 + 2x^3 + 3x^4,\quad (14 \text{ FLOPs})
\end{equation*}
we can write this equivalently as
\begin{equation*}
  p(x) = -5 + x(4 + x(-7 + x(2 + 3x))),\quad (8\text{ FLOPs})
\end{equation*}

\item This technique is called \textbf{nesting polynomials} or \textbf{Horner's Method}
\method{Nesting}{Standard polynomial}{Nested polynomial}{Factor out x to minimize the number of multiplications.}

\end{itemize}
\end{frame}
%%%%%%%%%%%%%%%%%%%%%%%%%%%%%%%%%%%%%%%%%%%
\begin{frame}{Nesting}

\begin{itemize}
    \item The standard polynomial costs $\mO(n^2)$ FLOPs to evaluate 
    \item The fully nested polynomial costs exactly $2n$ FLOPs
    \item Compare $$ p(x) = -5 + 4x -7x^2 + 2x^3 + 3x^4,$$ versus $$ p(x) = -5 + x(4 + x(-7 + x(2 + 3x)))$$
\end{itemize}

    \keyidea{Nested polynomials are evaluated more optimally.}
\end{frame}

%%%%%%%%%%%%%%%%%%%%%%%%%%%%%%%%%%%%%%%%%%%
\begin{frame}{Summary so far:}
\begin{itemize}
  
   \item \textbf{Monomial basis}: $p(x) = a_0 + a_1 x +\dots+a_n x^n$ results in a poorly conditioned Vandermonde matrix that can lead to inaccurate coefficients.\bigskip

  \item \textbf{Nesting}: It is trivial to minimize the cost of polynomial evaluation. \bigskip
  
  \item \textbf{Lagrange basis}: $p(x) = \ell_0(x) y_0 + \dots + \ell_n(x) y_n$. There is no error in the coefficients. It's easy to compute by hand.  \bigskip

  \begin{itemize}
    \item When adding an additional data point, you cannot re-use previous polynomial easily (for monomials and Lagrange)
  \end{itemize}

\end{itemize}
\end{frame}

%%%%%%%%%%%%%%%%%%%%%%%%%%%%%%%%%%%%%%%%%%%
\sectionframe{Newton Interpolation}
%%%%%%%%%%%%%%%%%%%%%%%%%%%%%%%%%%%%%%%%%%%
\begin{frame}{Monomial Basis}
\begin{itemize}
\item Consider interpolating the data
\begin{center}
  \begin{tabular}{c | c | c}
  $x_0$ & $x_1$ & $x_2$\\\hline
  $y_0$ & $y_1$ & $y_2$\\
  \end{tabular}
\end{center}
\item Our polynomial using a \textbf{monomial basis} will be   
\begin{equation*}
  p_2(x) = a + bx + cx^2
\end{equation*}

\item We can write out a system of equations for $a,b,c$
\[
\begin{alignedat}{4}
  \text{at }x_0:\quad 
    &p(x_0)=y_0= a &+ bx_0 &+ c x_0^2,\\
  \text{at }x_1:\quad 
    &p(x_1)=y_1 =a &+ bx_1 &+ c x_1^2, \\
  \text{at }x_2:\quad 
    &p(x_2)=y_2 =a &+ bx_2 &+ c x_2^2. \\
\end{alignedat}\implies
    \begin{bmatrix}
      1 & x_0 & x_0^2 \\
      1 & x_1 & x_1^2 \\
      1 & x_2 & x_2^2 
    \end{bmatrix}
    \begin{bmatrix}
      a\\
      b\\
      c\\
    \end{bmatrix}
     =
    \begin{bmatrix}
      y_0\\
      y_1\\
      y_2\\
    \end{bmatrix}
\]
\item Solve the $3\times3$ \textbf{Vandermonde} system
\end{itemize}
\end{frame}
%%%%%%%%%%%%%%%%%%%%%%%%%%%%%%%%%%%%%%%%%%%
\begin{frame}{Another Basis?}

\question{
What if we use this basis instead? 
\begin{center}
  \begin{tabular}{l | c}
  Function & Order \\\hline
  $1$ & 0\\
  $x-x_0$ & 1\\
  $(x-x_0)(x-x_1)$ & 2
  \end{tabular}
\end{center}
}
\begin{itemize}
    \item Then our interpolation would be 
    $$ p_2(x) = a + b(x-x_0) + c(x-x_0)(x-x_1).$$

    \item The three unknowns are $a$, $b$, $c$
\end{itemize}
\end{frame}
%%%%%%%%%%%%%%%%%%%%%%%%%%%%%%%%%%%%%%%%%%%
\begin{frame}{Conditions}

        \begin{itemize}
            \item The polynomial must pass through $(x_0,y_0)$, $(x_1,y_1)$, and $(x_2,y_2)$

            \item The conditions are 
            \begin{align*}
                y_0=p_2(x_0)\\ 
                y_1=p_2(x_1)\\ 
                y_1=p_2(x_2) 
            \end{align*}
    \item If $p_2(x) = a + b(x-x_0) + c(x-x_0)(x-x_1)$ then
    $$
    \begin{alignedat}{4}
  \text{at }x_0:\quad 
    &y_0  &={}&\,a
            &{}+&\,0
            &{}+&\,0,\\
  \text{at }x_1:\quad 
    &y_1  &={}&\,a
            &{}+&\,b\,(x_1 - x_0)
            &{}+&\,0,\\
  \text{at }x_2:\quad 
    &y_2  &={}&\,a
            &{}+&\,b\,(x_2 - x_0)
            &{}+&\,c\,(x_2 - x_0)(x_2 - x_1).
\end{alignedat}
    $$
        \end{itemize}
\end{frame}
%%%%%%%%%%%%%%%%%%%%%%%%%%%%%%%%%%%%%%%%%%%
\begin{frame}{Matrix Form}
\begin{itemize}
    \item The matrix form of 
    $$
    \begin{alignedat}{4}
    \text{at }x_0:\quad 
    &y_0  &={}&\,a
    &{}+&\,0
    &{}+&\,0,\\
    \text{at }x_1:\quad 
    &y_1  &={}&\,a
    &{}+&\,b\,(x_1 - x_0)
    &{}+&\,0,\\
    \text{at }x_2:\quad 
    &y_2  &={}&\,a
    &{}+&\,b\,(x_2 - x_0)
    &{}+&\,c\,(x_2 - x_0)(x_2 - x_1).
    \end{alignedat}
    $$
    is
\begin{equation*}
  \begin{bmatrix}
    1 & 0 & 0\\
    1 & x_1 - x_0 & 0\\
    1 & x_2 - x_0 & (x_2-x_0)(x_2-x_1)\\
  \end{bmatrix}
  \begin{bmatrix}
    a\\ b\\ c\\
  \end{bmatrix} =
  \begin{bmatrix}
    y_0\\ y_1\\ y_2\\
  \end{bmatrix}
\end{equation*}

% \item This is preferable to solving the \textbf{Vandermonde matrix}
% $$
%     \begin{bmatrix}
%       1 & x_0 & x_0^2 \\
%       1 & x_1 & x_1^2 \\
%       1 & x_2 & x_2^2 
%     \end{bmatrix}
%     \begin{bmatrix}
%       a\\
%       b\\
%       c\\
%     \end{bmatrix}
%      =
%     \begin{bmatrix}
%       y_0\\
%       y_1\\
%       y_2\\
%     \end{bmatrix}
% $$
\end{itemize}

\keyidea{
This is a \textbf{triangular system}! This is great news!
}
\end{frame}
%%%%%%%%%%%%%%%%%%%%%%%%%%%%%%%%%%%%%%%%%%%
\begin{frame}{Comparison}

\begin{block}{Monomial Basis}
 Fitting a polynomial of the form   $p(x) = a + bx + cx^2$  
 results in a \textbf{full matrix}
$$
    \begin{bmatrix}
      1 & x_0 & x_0^2 \\
      1 & x_1 & x_1^2 \\
      1 & x_2 & x_2^2 
    \end{bmatrix}
    \begin{bmatrix}
      a\\
      b\\
      c\\
    \end{bmatrix}
     =
    \begin{bmatrix}
      y_0\\
      y_1\\
      y_2\\
    \end{bmatrix}
$$ 
Costs $\mO(n^3)$ to solve
\end{block}


\begin{block}{Alternate Basis}
     Fitting a polynomial of the form   $ p = a + b(x-x_0) + c(x-x_0)(x-x_1)$  
 results in a \textbf{triangular matrix}
\begin{equation*}
  \begin{bmatrix}
    1 & 0 & 0\\
    1 & x_1 - x_0 & 0\\
    1 & x_2 - x_0 & (x_2-x_0)(x_2-x_1)\\
  \end{bmatrix}
  \begin{bmatrix}
    a\\ b\\ c\\
  \end{bmatrix} =
  \begin{bmatrix}
    y_0\\ y_1\\ y_2\\
  \end{bmatrix}
\end{equation*}
Costs $\mO(n^2)$ to solve
\end{block}


\end{frame}
%%%%%%%%%%%%%%%%%%%%%%%%%%%%%%%%%%%%%%%%%%%
\begin{frame}{Newton Polynomials}
\begin{itemize}
\item \textbf{Newton Polynomials} are of the form
  \begin{equation*}
    p_n(x) = a_0 + a_1(x-x_0) + a_2(x-x_0)(x-x_1) + a_3(x-x_0)(x-x_1)(x-x_2) + \dots
  \end{equation*}

\item The \textbf{Newton basis functions} have the form:
\begin{center}
  \begin{tabular}{l | c}
  Function & Order \\
  \hline
  $1$ & 0\\
  $x-x_0$ & 1\\
  $(x-x_0)(x-x_1)$ & 2\\
  $(x-x_0)(x-x_1)(x-x_2)$ & 3\\
  \end{tabular}
\end{center}
\vspace{3mm}
\define{Newton Basis}{

The \textbf{Newton basis functions } are the family of polynomials 
    $$\mathcal{F} = \{ 1, x-x_0, (x-x_0)(x-x_1), \ldots  \} $$
where $x_k$ are distinct data points.
}
\end{itemize}
\end{frame}
%%%%%%%%%%%%%%%%%%%%%%%%%%%%%%%%%%%%%%%%%%%%%%%%%%%%%%%%%%%%%%%%%%%%%%%%
\begin{frame}{Newton Interpolation Method}
    \method{Newton Interpolation}{Data $n+1$ points $(x_i,y_i)$ for $i=0,\ldots,n$}{Interpolating polynomial $p_n(x)$}{The system obtained with the Newton interpolation is triangular.}

\begin{itemize}
    \item Monomial vs. Newton 
    $$
        \begin{bmatrix}
      1 & x_0 & x_0^2 \\
      1 & x_1 & x_1^2 \\
      1 & x_2 & x_2^2 
    \end{bmatrix}
    \;\text{vs.}\;
      \begin{bmatrix}
    1 & 0 & 0\\
    1 & x_1 - x_0 & 0\\
    1 & x_2 - x_0 & (x_2-x_0)(x_2-x_1)\\
  \end{bmatrix}
    $$
\item Can be solved in $\mO(n^2)$ FLOPs 
\end{itemize}
\end{frame}

%%%%%%%%%%%%%%%%%%%%%%%%%%%%%%%%%%%%%%%%%%%%%%%%%%%%%%%%%%%%%%%%%%%%
\sectionframe{Solving for the Coefficients}
%%%%%%%%%%%%%%%%%%%%%%%%%%%%%%%%%%%%%%%%%%%%%%%%%%%%%%%%%%%%%%%%%%%%

\begin{frame}{Forward Substitution}
\begin{itemize}
    \item We can solve this system with forward substitution 
    \begin{equation*}
  \begin{bmatrix}
    1 & 0 & 0\\
    1 & x_1 - x_0 & 0\\
    1 & x_2 - x_0 & (x_2-x_0)(x_2-x_1)\\
  \end{bmatrix}
  \begin{bmatrix}
    a_0\\ a_1\\ a_2\\
  \end{bmatrix} =
  \begin{bmatrix}
    y_0\\ y_1\\ y_2\\
  \end{bmatrix}
\end{equation*}

\item First row:
    $$
    a_0 = y_0
    $$
\item Second row: 
\begin{align*}    
a_0 + (x_1-x_0) a_1 &= y_1\\ 
 a_1 &= \frac{y_1-a_0}{x_1-x_0}\\ 
\end{align*}
\item Third row:
\begin{align*}    
a_0 + (x_2-x_0) a_1 + (x_2-x_0)(x_2-x_1)a_2&= y_2\\ 
 a_2 &= \frac{y_2-a_0 - (x_2-x_0) a_1}{(x_2-x_0)(x_2-x_1)}\\ 
\end{align*}
\end{itemize}
\end{frame}

%%%%%%%%%%%%%%%%%%%%%%%%%%%%%%%%%%%%%%%%%%%%%%%%%%%%%%%%%%%%%%%%%%%%
\begin{frame}{Pattern for Coefficients?}
\question{Do you see a pattern?}
\begin{align*}
  a_0 & = y_0 = f(x_0)\\
  a_1 & = \frac{y_1 - a_0}{x_1 - x_0} = \frac{f(x_1)-f(x_0)}{x_1 - x_0}\\
  a_2 & = \frac{y_2 - a_0 - (x_2-x_0)a_1}{(x_2-x_1)(x_2-x_0)}\\
      & = \frac{f(x_2)-f(x_0)-(x_2-x_0)\frac{f(x_1)-f(x_0)}{x_1-x_0}}{(x_2-x_1)(x_2-x_0)}\\
      & = \frac{\frac{f(x_2)-f(x_1)}{x_2-x_1} - \frac{f(x_1)-f(x_0)}{x_1-x_0} }{x_2-x_0}
\end{align*}
\begin{center}
    There is a \textbf{recursive pattern} of $\frac{\Delta f}{\Delta x}$.
\end{center}
\end{frame}
%%%%%%%%%%%%%%%%%%%%%%%%%%%%%%%%%%%%%%%%%%%%%%%%%%%%%%%%%%%%%%%%%%%%
\begin{frame}{Divided Differences}
\begin{itemize}
    \item  We see many instances of the fraction
\begin{equation*}
  \frac{f(x_j)-f(x_i)}{x_j-x_i}
\end{equation*}
\item They're called \textbf{divided differences} and are denoted:
\begin{equation*}
  f[x_i,x_j] := \frac{f(x_j)-f(x_i)}{x_j-x_i}
\end{equation*}

\item For example 
    $$  f[a,b] = \frac{f(a) - f(b)}{a - b}$$

\item Notice that $a$ and $b$ can be switched 
$$ f[b,a] = \frac{f(b)-f(a)}{b-a} = \frac{-(f(a) - f(b)))}{-(a-b)} = \frac{f(a) - f(b)}{a - b} = f[a,b]$$
\end{itemize}
\end{frame}
%%%%%%%%%%%%%%%%%%%%%%%%%%%%%%%%%%%%%%%%%%%%%%%%%%%%%%%%%%%%%%%%%%%%%%%%
\begin{frame}{Bracket Notation}
\begin{itemize}
    \item The \textbf{zeroth order divided difference} is just a base case 
      $$a_0= f[x_0] = f(x_0)$$ 
    \item The \textbf{first order divided difference} is defined to be
     $$ a_1 = f[x_1,x_0] = \frac{f[x_1]-f[x_0]}{x_1-x_0} $$
    \item The \textbf{second order divided difference} is 
     $$ a_2 = f[x_0,x_1,x_2] = \frac{f[x_1,x_2]-f[x_0,x_1]}{x_2-x_0} $$
    \item For higher order terms we use recursion
\end{itemize}
\begin{block}{Recursive Property}
$$ f[x_0,\dots,x_k] = \frac{f[x_1,\dots,x_k]-f[x_0,\dots,x_{k-1}]}{x_k - x_0}$$
\end{block}
\end{frame}
%%%%%%%%%%%%%%%%%%%%%%%%%%%%%%%%%%%%%%%%%%%%%%%%%%%%%%%%%%%%%%%%%
\begin{frame}{Recursive Property}

\begin{block}{Recursive Property}
$$ f[x_0,\dots,x_k] = \frac{f[x_1,\dots,x_k]-f[x_0,\dots,x_{k-1}]}{x_k - x_0}$$
\end{block}

$$ f\underbrace{[x_0,\dots,x_k]}_{0 \text{ to } k} = \frac{\underbrace{[x_1,\dots,x_k]}_{1 \text{ to } k}-f\underbrace{[x_0,\dots,x_{k-1}]}_{0 \text{ to }k-1}}{x_k - x_0}$$

\begin{itemize}
    \item Let's call $f[x_0,\dots,x_k]$ the $k^{th}$ level divided difference
    \item $k^{th}$ level divided difference requires two $(k-1)^{th}$ level divided differences
\end{itemize}


\end{frame}
%%%%%%%%%%%%%%%%%%%%%%%%%%%%%%%%%%%%%%%%%%%%%%%%%%%%%%%%%%%%%%%%%

\begin{frame}{The Pattern}
\begin{itemize}
    \item These are the coefficients 
\begin{align*}
  a_0 &  = f(x_0)\\
  a_1 & = \frac{f(x_1)-f(x_0)}{x_1 - x_0}\\
  a_2 & = \frac{\frac{f(x_2)-f(x_1)}{x_2-x_1} - \frac{f(x_1)-f(x_0)}{x_1-x_0} }{x_2-x_0}
\end{align*}

\item They can be expressed compactly using the \textbf{bracket notation}
\begin{align*}
  a_0 &  = f[x_0]\\
  a_1 & = f[x_0,x_1]\\
  a_2 & = \frac{f[x_1, x_2] - f[x_0, x_1] }{x_2-x_0} =: f[x_0, x_1, x_2]
\end{align*}

\end{itemize}
\end{frame}
%%%%%%%%%%%%%%%%%%%%%%%%%%%%%%%%%%%%%%%%%%%%%%%%%%%%%%%%%%%%%%%%%%%
\begin{frame}{Newton Polynomial Coefficients}
\begin{itemize}
    \item The Newton polynomial is 
    $$
    p(x) = a_0 + a_1(x-x_0) + a_2(x-x_0)(x-x_1)+a_3(x-x_0)(x-x_1)(x-x_2) + \text{etc.}
    $$
    \item The coefficients are
    \begin{align*}
  a_0 &  = f[x_0]\\
  a_1 & = f[x_0,x_1]\\
  a_2 & =  f[x_0, x_1, x_2]\\
  a_3 & =  f[x_0, x_1, x_2,x_3]\\
  & \vdots\\ 
  a_n &= f[x_0,x_1\dots ,x_n]
\end{align*}
\end{itemize}
\keyidea{Just add one extra point in the bracket. For example
$$
a_{n+1} = f[x_0,x_1,\dots,x_n,x_{n+1}]
$$}
\end{frame}

%%%%%%%%%%%%%%%%%%%%%%%%%%%%%%%%%%%%%%%%%%%%%%%%%%%%%%%%%%%%%%%%%%%
\begin{frame}[shrink]{Newton Polynomials using Divided Differences}
\begin{example}
  Find the 2nd order interpolating Newton polynomial for the data points 
  \begin{center}
    \begin{tabular}{c | c c c}
      $x$ & 1 & -4 & 0\\\hline
      $y$ & 3 & 13 & -23\\
    \end{tabular}
  \end{center}
\end{example}
\begin{itemize}
    \item[1.] First, write the general 2nd degree Newton polynomial
\begin{equation*}
  p_2(x) = a_0  + a_1(x-x_0) + a_2(x-x_0)(x-x_1)
\end{equation*}
\item[2.] Next, compute the coefficients
\begin{align*}
  a_0 & = f[x_0] = f[1] = f(1) = 3\\
  a_1 & = f[x_0,x_1] = \frac{f(x_1)-f(x_0)}{x_1 - x_0} = \frac{13-3}{-4-1} =
  -2\\
  a_2 & = f[x_0,x_1,x_2] = \frac{f[x_1,x_2]-f[x_0,x_1]}{x_2-x_0} = \frac{\frac{-23-13}{0--4} - \frac{13-3}{-4-1}}{0-1} = \frac{-9+2}{-1} = 7
\end{align*}
\item[3.] Finally, substitute the values
\begin{equation*}
  p_2(x) = 3 -2(x-1) + 7(x-1)(x+4)
\end{equation*}
\end{itemize}
\end{frame}
%%%%%%%%%%%%%%%%%%%%%%%%%%%%%%%%%%%%%%%%%%%%%%%%%%%%%%%%%%%%%%%
% \begin{frame}{Nesting}
% \begin{itemize}
%     \item The Newton polynomial $p_2(x) = 3 -2(x-1) + 7(x-1)(x+4)$ costs 2 additions and 3 additions = \textbf{5 FLOPs}
% \vfill 
%     \item The standard form $p_2(x)=7 x^2 + 19 x - 23$ would cost \textbf{6 FLOPs}
% \vfill 
%     \item The most optimal form costs \textbf{4 FLOPs} 
% \vfill 
% \end{itemize}
%     \keyidea{The Newton polynomial naturally introduces partial nesting.}
% \end{frame}
%%%%%%%%%%%%%%%%%%%%%%%%%%%%%%%%%%%%%%%%%%%%%%%%%%%%%%%%%%%%%%%
\begin{frame}{Invariance}
\begin{block}{Invariance Theorem}
    $f[x_0,\dots,x_k]$ is invariant under all permutations of the arguments
    $x_0,\dots,x_k$
\end{block}
\vspace{4mm}

Why is this important? 
\begin{itemize}
    \item The nodes $x_i$ do not have to arranged from least to greatest 
    \item We can apply the recursive formula easily 
    \item If an additional interpolation point is added $(x_{n+1}, y_{n+1})$, we can easily find the next coefficient
\end{itemize}


\end{frame}
%%%%%%%%%%%%%%%%%%%%%%%%%%%%%%%%%%%%%%%%%%%%%%%%%%%%%%%%%%%%%%%%%%%%%
\begin{frame}{Example}
    \begin{block}{Existing Interpolation}
      \begin{center}
    \begin{tabular}{c | c c c}
      $x$ & 1 & -4 & 0\\\hline
      $y$ & 3 & 13 & -23\\
    \end{tabular}
  \quad $p_2(x) = 3 -2(x-1) + 7(x-1)(x+4)$
  \end{center}
    \end{block}
    \begin{block}{Extra Point}
        If we add an extra point
              \begin{center}
    \begin{tabular}{c | c c c c}
      $x$ & 1 & -4 & 0   &2\\\hline
      $y$ & 3 & 13 & -23 &3 \\
    \end{tabular}
  \end{center}
   then our new interpolation is
   $$p_3(x) = 3 -2(x-1) + 7(x-1)(x+4)+a_3(x-1)(x+4)(x)$$
    \end{block}

    \keyidea{
    When adding a new point, we may reuse the previously calculated coefficients.
    }
\end{frame}


%%%%%%%%%%%%%%%%%%%%%%%%%%%%%%%%%%%%%%%%%%%%%%%%%%%%%%%%%%%%%%%%%%%%%
\begin{frame}{Newton Interpolation}
    \method{Newton Divided Differences}{$n+1$ data points $(x_i,y_i)$}{Degree $n$ interpolating polynomial $p_n(x)$}{Take advantage of the recursive formula of divided differences}

    \keyidea{
    Advantages to using Newton interpolation:
\begin{itemize}
        \item Newton system can be solved in $\mO(n^2)$ FLOPs
        % \item Newton polynomials are partially nested
        \item When an additional points, we can build off of the existing interpolation
        \item<2> Can be easily calculated by hand using the \textbf{tabular method}
    \end{itemize}
    }
\end{frame}

%%%%%%%%%%%%%%%%%%%%%%%%%%%%%%%%%%%%%%%%%%%%%%%%%%%%%%%%%%%%%%%%%%%%%%%%
\sectionframe{Tabular Method}
%%%%%%%%%%%%%%%%%%%%%%%%%%%%%%%%%%%%%%%%%%%%%%%%%%%%%%%%%%%%%%%%%%%%%%%%
\begin{frame}{Divided Differences with the Table}
\begin{block}{Recursive Property}
Use this to fill out the table
$$ f[x_0,\dots,x_k] = \frac{f[x_1,\dots,x_k]-f[x_0,\dots,x_{k-1}]}{x_k - x_0}$$
\end{block}
We can compute the divided differences much easier using tables.  To
construct the divided difference table for $f(x)$ for the $x_0,\dots,x_3$
\begin{center}
\begin{tabular}{l | l | l | l | l}
$x$ & $f[\cdot]$ & $f[\cdot,\cdot]$ & $f[\cdot,\cdot,\cdot]$ & $f[\cdot,\cdot,\cdot,\cdot]$\\\hline
\alt<2,5,7>{\color{red}$x_0$}{$x_0$}    & \alt<2>{\color{red}$f[x_0]$}{$f[x_0]$}&              &                  &                      \\
                                     &                                        & \alt<2,5>{\color{red}$f[x_0,x_1]$}{$f[x_0,x_1]$} &                  &                      \\
\alt<2,3,6>{\color{red}$x_1$}{$x_1$}  & \alt<2,3>{\color{red}$f[x_1]$}{$f[x_1]$} &              & \alt<5,7>{\color{red}$f[x_0,x_1,x_2]$}{$f[x_0,x_1,x_2]$} &                      \\ 
                                     &                                        & \alt<3,5,6>{\color{red}$f[x_1,x_2]$}{$f[x_1,x_2]$} &                  & \alt<7>{\color{red}$f[x_0,x_1,x_2,x_3]$}{$f[x_0,x_1,x_2,x_3]$} \\
\alt<3,4,5>{\color{red}$x_2$}{$x_2$}  & \alt<3,4>{\color{red}$f[x_2]$}{$f[x_2]$}   &              & \alt<6,7>{\color{red}$f[x_1,x_2,x_3]$}{$f[x_1,x_2,x_3]$} &                      \\
                                     &                                        & \alt<4,6>{\color{red}$f[x_2,x_3]$}{$f[x_2,x_3]$} &                  &                      \\
\alt<4,6,7>{\color{red}$x_3$}{$x_3$}  & \alt<4>{\color{red}$f[x_3]$}{$f[x_3]$} &              &                  &                      \\
\end{tabular}
\end{center}
\end{frame}
%%%%%%%%%%%%%%%%%%%%%%%%%%%%%%%%%%%%%%%%%%%%%%%%%%%%%%%%%%%%%%%%%%%%%%%%

\begin{frame}{Tabular Method Example}
\begin{example}
Construct the divided differences table for the data
\begin{center}
\begin{tabular}{c | c c c c}
 \smallskip 
 $x$ & 1 & 3 & 0 & 2\\
 \hline
 $y$ & 0 & 6 & 3 & 7\\
\end{tabular}
\end{center}
and construct the largest order interpolating polynomial.\\\smallskip
\end{example}

Start by filling the first 2 columns with the given data. Then fill each column from left to right. 

\begin{center}
\begin{tabular}{c | c | c | c | c}
$x$ & $f[\cdot]$ & $f[\cdot,\cdot]$ & $f[\cdot,\cdot,\cdot]$ & $f[\cdot,\cdot,\cdot,\cdot]$\\\hline
\alt<2,5,7>{\color{red}$1$}{$1$}  & \alt<2>{\color{red}$0$}{$0$} &              &                  &                      \\
                                  &                               & \onslide<2->{\alt<2,5>{\color{red}$3$}{$3$}} &                  &                      \\
\alt<2,3,6>{\color{red}$3$}{$3$}  & \alt<2,3>{\color{red}$6$}{$6$} &            & \onslide<5->{\alt<5,7>{\color{red}$2$}{$2$}} &  \\ 
                                  &                                & \onslide<3->{\alt<3,5,6>{\color{red}$1$}{$1$}} &                  & \onslide<7->{\alt<7>{\color{red}$-3$}{$-3$}} \\
\alt<3,4,5>{\color{red}$0$}{$0$}  & \alt<3,4>{\color{red}$3$}{$3$}  &           & \onslide<6->{\alt<6,7>{\color{red}$-1$}{$-1$}} &                      \\
                                  &                                 & \onslide<4->{\alt<4,6>{\color{red}$2$}{$2$}} &                  &                      \\
\alt<4,6,7>{\color{red}$2$}{$2$}      & \alt<4>{\color{red}$7$}{$7$} &          &                  &                      \\
\end{tabular}
\end{center}
\end{frame}

%%%%%%%%%%%%%%%%%%%%%%%%%%%%%%%%%%%%%%%%%%%%%%%%%%%%%%%%%%%%%%%%%%%%%%%%
\begin{frame}{Writing the polynomial}
\begin{center}
\begin{tabular}{c | c | c | c | c}
$x$ & $f[\cdot]$ & $f[\cdot,\cdot]$ & $f[\cdot,\cdot,\cdot]$ & $f[\cdot,\cdot,\cdot,\cdot]$\\\hline
$1$  & {\color{red}$0$} &    &                  &                      \\
      &             & {\color{red}$3$}&                  &                      \\
{$3$}  &  {$6$} &    &  {\color{red}$2$} &  \\ 
      &               & {$1$}&                  &  {\color{red}$-3$} \\
{$0$}  & {$3$}  &           & {$-1$}&                      \\
       &                                 & {$2$} &                  &                      \\
{$2$}   & {$7$} &          &                  &                      \\
\end{tabular}
\end{center}
\begin{itemize}
\item After completing the table, write the general form for a cubic
$$
p_3(x) = a_0 + a_1(x-x_0) + a_2(x-x_0)(x-x_1) -  a_3 (x-x_0)(x-x_1)(x-x_2)
$$      
\item The coefficients $a_0$ through $a_3$ is along the upper diagonal
\begin{align*}
  p_2(x) &= {\color{red}{0}} + {\color{red}{3}}(x-1) + {\color{red}{2}}(x-1)(x-3)   {\color{red}{-3}}(x-1)(x-3)(x-0)\\
\end{align*}

\end{itemize}
\end{frame}

%%%%%%%%%%%%%%%%%%%%%%%%%%%%%%%%%%%%%%%%%%%%%%%%%%%%%%%%%%%%%%%%%%%%%%%%%%%%%%%
\begin{frame}{Math Joke Revisited}

What is the next number in the sequence? 
$$ 1, 2, 3,\underline{\hspace{ 6 mm}} $$

\onslide<2->{
It's whatever you want it to be!

\vspace{3mm}

Create an interpolating polynomial that passes through 
$$ (1,1), (2,2), (3,3), (4,x)$$ 
}

\end{frame}

\end{document}
